\subsection{What the split criterion sees, and what it discards}
\label{sec:variance}

Two properties of variance-based splitting underlie the effects in the sections
that follow. Take a
noise-free, sorted dataset and sweep every single split, scoring each by the sum
of within-child variances, for $y = x$, $y = x^2$, $y = x^3$ ($x =
1\ldots1000$):

\begin{center}
\begin{tabular}{cccc}
\toprule
response & min split-variance & fraction of parent variance removed & argmin split index (1--1000) \\
\midrule
$x$   & $4.2 \times 10^{4}$  & $0.75$ & 500 (centre) \\
$x^2$ & $4.3 \times 10^{10}$ & $0.76$ & 697 \\
$x^3$ & $4.1 \times 10^{16}$ & $0.74$ & 799 \\
\bottomrule
\end{tabular}
\end{center}

Here \code{argmin} is the position among the $1000$ sorted values of $x$ at
which the split minimizes the sum of within-child variances.

The raw minimum grows by orders of magnitude with the degree, but that is a
\emph{response-scale artifact}: as a fraction of parent variance removed, all
three splits are essentially equal (${\approx}0.75$). The two relevant
properties that the table illustrates are:

\begin{enumerate}[leftmargin=1.6em]
\item \textbf{Trees discard feature-distribution shape.} A split compares only
the
  \emph{ordering} of a feature's values, so the tree is invariant to any
  monotone
  transform and cannot see whether a feature is uniform or heavy-tailed,
  information a per-unit effect size needs. This is the root cause of the
  confound
  in Section~\ref{sec:confound}.
\item \textbf{The optimal split migrates with response shape.} For a linear
  response the best split is central; for higher orders it moves toward the
  high-$y$ tail ($500 \rightarrow 697 \rightarrow 799$). The imbalance leaves
  more
  structure in the large child, so more splits follow down that branch, each
  charged to the same feature.
\end{enumerate}

Neither property depends on \code{mtry}: the sweep has a single feature, and
\code{mtry} decides which features compete at a node, not where a feature's
best threshold falls or what the split can see.
Impurity importance is a sum of those split scores over the fitted trees, so it
inherits both properties: it summarises variance contribution over the feature
ordering, but is blind to the metric the ordering lives on.
